% ---- From Santi (removed by Carlos)
% Physics event generators are a critical component of simulations in high-energy physics, used to produce synthetic collision events for both theory studies and experimental analyses. As the field moves into the High Luminisity LHC (\hilumi) era, the demand for simulated events is skyrocketing~\cite{HSFPhysicsEventGeneratorWG:2020gxw}, while traditional CPU-based computing improvements have slowed. Projections indicate that flat computing budgets and the end of Moore's law speedups could leave future experiments limited by Monte Carlo (MC) statistics if generator performance is not dramatically improved. 

% Need to add references to LSZ formula and so on probably
Due to the intrinsically probabilistic nature of fundamental particle interactions and the composite structure of the proton, accurately modelling the dynamics of proton collisions at the Large Hadron Collider (LHC) constitutes one of the major challenges in high-energy physics (HEP). Physics event generators encode the necessary tools to describe these collisions in a programmatic way, enabling the use of computational resources to make quantitative predictions for a given theoretical model.

Each collision at the LHC is characterized by an initial and a final state. The particles in the initial state originate from the colliding protons, whereas the final-state particles are produced as a result of the interaction. It can be shown that the transition probability (\(\mathcal{P}\)) for an initial state \(i\), consisting of two incoming partons \(q_1\) and \(q_2\), to evolve into a final state \(\mathcal{X}\) is proportional to the squared matrix element \(\mathcal{M}\), as shown in Equation~\ref{eq:hard_scatter_bracket}.

\begin{equation}
    \mathcal{P}(i \rightarrow \mathcal{X}) = |\mathcal{M}_{i \rightarrow \mathcal{X}}|^2
    \label{eq:hard_scatter_bracket}
\end{equation}

To compute the probability of a given process \(\mathcal{X}\) occurring at the LHC, one must integrate the squared matrix element over the full phase space that connects the initial and final states. These calculations can only be performed up to a fixed order in a perturbation theory that expands over different powers of the so-called coupling constant, that quantifies the strength of the interaction between particles. To do so, physics event generators rely on Monte Carlo (MC) integration techniques, which approximate the integral by sampling random points that populate the accessible phase space. All modern MC generators provide calculations at leading order (LO), while some are capable of higher-order predictions such as next-to-LO (NLO) and next-to-NLO (NNLO). These higher order corrections are not negligible in most of the cases, and with the ever increasingly size of the LHC dataset, experiments at the LHC have become sensitive to these corrections. However, increasing the perturbative order leads to a rapid, non-linear growth in computational complexity, placing severe demands on computing resources.

% ---- From Santi (removed by Carlos)
% In complex LHC processes, the parton-level event generation stage can consume a large fraction of the overall simulation time – for example, the calculation of matrix elements (the squared quantum amplitudes for a given process) can account for on the order of 30–40\% of the total event generation CPU time in multi-jet processes


% This is because for each simulated event, the generator must evaluate a complicated physics matrix element function at a random point in phase space, integrating over many Feynman diagram contributions. Generating millions of events thus entails repeating these calculations millions of times. In principle, this task is embarrassingly parallel: the same matrix element evaluation is performed independently for many different phase-space points. This parallelism has motivated efforts to accelerate event generators using modern hardware architectures. Notably, GPUs have been applied to MadGraph5\_aMC@NLO~\cite{Alwall:2014hca}, yielding substantial speedups in matrix element computation (often by one to two orders of magnitude for complex processes) by exploiting their massive parallel throughput~\cite{Valassi:2021ljk,Valassi:2025xfn,Hageboeck:2023blb,Valassi:2023yud}.

As a result, as the volume of data recorded by the LHC experiments continues to grow, so do the requirements for accurate MC simulations. To address this growing computational burden, developers of MC generators have devoted significant effort to parallelizing matrix element integration using vectorized CPU and GPU infrastructures.

While GPUs and CPUs have thus demonstrated the value of heterogeneous computing for event generation~\cite{Carrazza:2021gpx}, the use of Field-Programmable Gate Arrays (FPGAs) remains relatively unexplored in this domain~\cite{Barbone:2023mul}. FPGAs allow custom hardware architectures tailored to the specific computation, offering the potential for fine-grained parallelism and pipeline concurrency that can drastically speed up repetitive arithmetic-heavy tasks. Importantly, FPGAs often operate at lower clock frequencies but with much lower overhead, translating to excellent energy efficiency. For certain workloads, FPGAs have been shown to be significantly more energy-efficient than GPUs – in some cases over an order of magnitude in performance-per-watt. This makes FPGA acceleration an attractive approach to sustainable high-throughput computing for HEP, where power and cooling constraints are growing concerns. However, historically FPGAs have been challenging to program for complex algorithms like those in MC generators, which involve hundreds of operations and complex control flow generated from Feynman rules. Recent advances in high-level synthesis (HLS) tools and heterogeneous devices aim to lower these barriers. In particular, the AMD Versal Adaptive Compute Acceleration Platform (ACAP) is a new FPGA-based system-on-chip that integrates not only reconfigurable logic but also an array of specialized AI Engine cores and that promises to deliver exceptional compute density with low latency and power usage. The Versal platform thus provides a versatile heterogeneous playground to map both custom logic and parallel DSP workloads, and it has been demonstrated to deliver ``high throughput, low latency, and power efficient'' processing for suitable applications. 